\section{Adaptive machines: trace classes and the popular trace}\label{sec:traces}

The cut lemmas of \cref{sec:locality,sec:blocks} assume obliviousness: all inputs share one head
trajectory, so the entry data of a step set is a well-defined finite object. A general machine
moves its heads depending on the data. The way around this is to group the inputs by their
\emph{block-level trace}, a coarse record of the head motion that costs only $O(db)=o(L)$ bits when the running time is $O(L)$.
Some trace is shared by a $2^{-O(db)}$ fraction of all inputs (the \emph{popular trace}), and on
that set the machine behaves like an oblivious one at the level of blocks. The cut lemmas then hold
with an additive loss of $O(db)$ bits per label class.

\subsection{Block-level traces}

Let $r\ge1$ be the block size.

\begin{definition}[Block-level trace, trace class]\label{def:trace}
\formal{Tape/DefsTrace.lean}{blockTrace, SameTrace, TraceClass}
The \emph{block-level trace} of the run of $M$ on $w$ at slot $y$ consists of the pair residency of
every work tape at slot $y$, the set of input rows touched during slot $y$, and the output length
at the first step $yr$ of the slot. Two inputs have the \emph{same trace} (for a parameter $T$) if
their traces agree at all slots $y\le\lceil T/r\rceil$. A \emph{trace class} with representative
$w_0$ is a set $F$ of matrix inputs containing $w_0$, all of which halt after exactly $T$ steps and
have the trace of $w_0$.
\end{definition}

\begin{lemma}[Trace classes generalize oblivious sets]\label{lem:trace-obl}
\formal{Tape/Statements/Adaptive.lean}{traceClass_of_oblivious, traceClass_of_obliviousOn, traceClass_of_mem}
If $M$ is oblivious on $\mathrm{Mat}(d,b)$ with halting time $T$, all matrix inputs form one trace
class, of size $2^{d^2b}$. More generally, every nonempty set of matrix inputs on which $M$ is oblivious (with halting time $T$) is a trace class, with any of its members as representative; and any member of a trace class is a
representative.
\end{lemma}

\begin{outline}
The trace is a function of the head trajectories and the output schedule.
\end{outline}

\subsection{Block locality}

For a trace class the entry data of a union of slots must record, at each memory boundary, more than
the offsets of the heads within their blocks: two inputs with the same trace have the same resident blocks,
but a head need not stand in the same one of them in both runs (\cref{prop:offsets}). We record the full
head positions.

\begin{lemma}[Block Locality]\label{lem:block-locality}
\formal{Tape/Statements/Adaptive.lean}{block_locality}
Let $r\ge1$ and let $w,w'$ be inputs of length $d^2b$ with the same trace. Let $U$ be a set of
slots, and suppose that
\begin{enumerate}[label=(\roman*)]
\item at the first step $zr$ of every memory boundary $z$ of $U$ the two runs agree on the state,
the positions of all heads, and the contents of every cell of every block resident at $z$;
\item for every chord of the pair residency of $w$ entering $U$ at a slot $z$ for a block $\beta$,
the two runs agree at time $zr$ on the contents of the cells of $\beta$; and
\item $w$ and $w'$ agree on every input bit in a row touched during $U$.
\end{enumerate}
Then at every step of the slots of $U$ the two runs agree on the state, all head positions, the
output length, the symbols under the work heads and the transition taken.
\end{lemma}

No halting time and no correctness are assumed.

\begin{outline}
This is the entry form of the Locality Lemma (\lean{locality_entry}, a corollary of
\cref{lem:locality-core}), with step $0$ and the first steps of the memory boundaries as anchors.
For a read of a cell in block $\beta$ at a step $\tau$ of the slots of $U$, the charging argument
of \cref{lem:charging} provides the time $a$: either $a=0$, when $\beta$ is resident at no slot
outside $U$ before $\tau$'s slot, or $a=zr$, where $z$ is the first slot after the last such slot
at which $\beta$ is resident, a memory boundary or the target of a chord entering $U$. Every visit
to the cell during $[a,\tau)$, in either run, lies in $U$, because the head's block is always
resident and the residencies of the two runs coincide.
\end{outline}

Recording only the offsets of the work heads within their blocks does not suffice:

\begin{proposition}[Offsets are not enough]\label{prop:offsets}
\formal{Tape/Statements/Adaptive.lean}{block_locality_offsets_counterexample}
There are a machine $M$ that transposes $1\times1$ matrices with $1$-bit entries and two inputs
$w,w'$, both halting at time $9$ with the same trace for $r=4$, such that slot $2$ is a memory
boundary of $U=\{2\}$, the runs agree at time $8$ on the state, the input-head position, the
work-head offsets modulo $4$ and the contents of the resident blocks, no chord enters slot $2$, and
no input row is touched during slot $2$; yet the work heads are in different blocks at time $8$
and the transitions at step $8$ differ.
\end{proposition}

\begin{outline}
An explicit machine, evaluated on both inputs.
\end{outline}

\subsection{The Cut Lemma on trace classes}

Write $G'=\clog{(g+1)}$, and let $D$ be such that $2^{d^2b}\le 2^D\abs{F}$ (the co-entropy of the
class).

\begin{lemma}[Cut Lemma on a trace class]\label{lem:cut-trace}
\formal{Tape/Statements/Adaptive.lean}{cut_bound_trace, cut_bound_trace_bits, cut_bound_trace_on, cut_bound_trace_bits_on}
Let $r\ge1$, let $F$ be a trace class with representative $w_0$, and let $M$ be correct on $F$ (for
instance, let $M$ transpose). For every set $U$ of slots, with $\mu=\mu(U)$ and $\chi=\chi(U)$ taken
for the pair residency of the representative,
\[
  2^{b\abs{S_{J(U)}}}\,\abs{F}\;\le\;2^{d^2b}\Bigl((q+1)(2T+1)^{k+1}\Bigr)^{\mu}(g+1)^{r(3k\mu+\chi)},
\]
\[
  b\,\abs{S_{J(U)}}\;\le\;D+\mu\bigl(\clog{(q+1)}+(k+1)\clog{(2T+1)}\bigr)+r\,G'\,(3k\mu+\chi).
\]
\end{lemma}

\begin{outline}
Inject the class into pairs (entry data, input bits outside a row-based cut set) through the counting core
of \cref{lem:cut-on}: since members read different positions but the same rows, the cut set is replaced by
the bits of the rows not touched during $J(U)$ whose transposed copies are appended during $J(U)$; it
contains the $b\abs{S_{J(U)}}$ bits of the cut pairs. The entry data consist of the state and the head
positions (which lie in $[-T,T]$) at each memory boundary, together with the contents of the resident blocks
there and of the blocks entering along chords. The rows touched and the output positions appended during
$U$ are class invariants, since the rows of each slot and the output lengths at slot starts are part of the
trace and all members halt at $T$. Injectivity follows from \cref{lem:block-locality} and correctness on
$F$.
\end{outline}

\begin{corollary}[Discrete Isoperimetry on a trace class]\label{cor:isoperimetry-trace}
\formal{Tape/Statements/Adaptive.lean}{isoperimetry_trace, isoperimetry_trace_on}
Under the hypotheses of \cref{lem:cut-trace}, with $1\le r\le db$ and a labelling of the slots with
$\Mb$ label changes and at most $m$ slots per label,
\[
  (d\dotminus\Mb)(d\dotminus2m)\,b\;\le\;(\Mb+1)D+\Mb\bigl(\clog{(q+1)}+(k+1)\clog{(2T+1)}\bigr)
  +r\,G'\,(3k\Mb+\chicross),
\]
where $x\dotminus y=\max(x-y,0)$ and $\chicross$ is taken for the representative's pair residency.
\end{corollary}

\begin{outline}
Sum the bits form of \cref{lem:cut-trace} over the at most $\Mb+1$ label classes and bound the left
side as in \cref{cor:isoperimetry}.
\end{outline}

\subsection{The popular trace}

\begin{definition}[Slot length and horizon]\label{def:slotlen}
\formal{Tape/DefsM.lean}{cellBits, slotLen, horizon}
For a machine $M$ let $G=\max(1,\clog{(g+1)})$ be the number of bits per work cell. The slot
length at size $d$ is $r=\lceil db(d)/G\rceil$, so that a block of $r$ cells holds about one matrix row.
For a time constant $c$ the \emph{horizon} is $c_0=cG+1$; if $T\le cL$ there are at most $c_0d$
slots (\lean{numSlots_le_horizon}).
\end{definition}

\begin{theorem}[Popular trace]\label{thm:popular}
\formal{Tape/Statements/Adaptive.lean}{exists_popular_traceClass, traceBits_le}
Let $1\le r\le db$ and suppose every matrix input halts within $T_{\max}$ steps. Then for some
$T\le T_{\max}$ there is a trace class $F$ with $2^{d^2b}\le\mathcal{N}\cdot\abs{F}$, where
$\mathcal{N}=\mathcal{N}(k,d,\lceil T_{\max}/r\rceil,T_{\max})$ is an explicit bound on the number
of pairs (halting time, trace) (\lean{traceCount}). If $T_{\max}\le cL$ and $r$ is the slot length,
then $\clog{\mathcal{N}}\le C\,d\,b(d)$ for an explicit constant $C$ depending only on $M$ and $c$
(\lean{traceConst}). So the popular class has co-entropy $D=O(db)=o(L)$.
\end{theorem}

\begin{outline}
Pigeonhole over the pairs (halting time, trace): block indices lie in $[-N,N]$, each residency has
at most three blocks, each slot touches at most two rows, and output lengths are at most
$T_{\max}$. The bound on $\clog{\mathcal{N}}$ is arithmetic from $T\le cd^2b$ and $b\le d\le2^b$.
\end{outline}

The trace used here records absolute block indices, rows and output lengths per slot, which costs
$O(d\log d)$ bits for running time $O(L)$. For machines that write their whole output at the end, a coarser
trace that records only block and row moves and residency patterns, $O(1)$ bits per slot, suffices, and the popular class then has co-entropy $O(d)$; \cref{sec:conjecture} uses this to show
that the coding conjecture is needed only at co-entropy $O(d)$.
